%\begin{longtable}{p{0.12\textwidth}p{0.22\textwidth}p{0.26\textwidth}p{0.25\textwidth}}
\setlength{\emergencystretch}{3em}
\begin{longtable}{
  p{0.12\textwidth}
  p{0.22\textwidth}
  p{0.26\textwidth}
  p{\dimexpr0.40\textwidth-8\tabcolsep\relax}
}
\caption{Preguntas del dataset, solucion aplicada e interpretacion.}

\label{tab:preguntas-consultas}\\

\toprule

\textbf{Item} & \textbf{Pregunta que responde} & \textbf{Columnas/\newline caracteristicas usadas} & \textbf{Solucion e \newline interpretacion} \\

\midrule

\endfirsthead

\toprule

\textbf{Item} & \textbf{Pregunta que responde} & \textbf{Columnas/\newline caracteristicas usadas} & \textbf{Solucion e \newline interpretacion} \\

\midrule

\endhead

ETL y tablas JSON &

¿Como separar el dataset original en estructuras reutilizables sin depender de PostgreSQL? &

\texttt{COD\_CONTRIBUYENTE}, \texttt{NOM\_CONTRIBUYENTE}, \texttt{COD\_PREDIO}, \texttt{PERIODO}, montos, descuentos y \texttt{UBIGEO}. &

Se generan \texttt{contribuyentes.json} y \texttt{emisiones.json}. La primera funciona como dimension y la segunda como tabla de hechos a nivel predio-contribuyente-periodo. Codigo: \texttt{ETL.scala}. \\

\midrule

SQL 1: \texttt{groupBy} + funciones &

¿Cuanto se emitio en total y en promedio por concepto municipal durante cada periodo? &

\texttt{PERIODO}, \texttt{MONTO\_PARQUE\_JARDIN}, \texttt{MONTO\_SERENAZGO}, \texttt{MONTO\_RESIDUO\_SOLIDO}, \texttt{MONTO\_BARRIDO\_CALLE}. &

Se agrupa por periodo y se aplican \texttt{sum}, \texttt{avg}, \texttt{count} y \texttt{round}. Permite comparar la evolucion anual de la emision por concepto. Codigo: \texttt{SqlQueries.query1GroupByPeriodo}. \\

\midrule

SQL 2: vista temporal + \texttt{ORDER BY} &

¿Cuales son los predios exclusivos con mayor monto total neto en 2025? &

\texttt{COD\_PREDIO}, \texttt{COD\_CONTRIBUYENTE}, \texttt{PERIODO}, \texttt{PORCENTAJE\_CONDOMINIO}, montos y descuentos. &

Se crea una vista temporal y se filtra 2025 con condominio 100. Luego se calcula monto neto y se ordena de mayor a menor. Codigo: \texttt{SqlQueries.query2TopPrediosTempView}. \\

\midrule

Regresion MLP &

¿Se puede estimar el monto de serenazgo a partir de variables monetarias y del periodo? &

\texttt{MONTO\_SERENAZGO} como etiqueta; \texttt{PERIODO}, \texttt{PORCENTAJE\_CONDOMINIO}, montos y descuentos como features. &

Se implementa un MLP de regresion con salida lineal y perdida MSE. El menor RMSE frente a SVR indica mejor ajuste global para esta tarea. Codigo: \texttt{MLPRegressor.scala}. \\

\midrule

Regresion SVR &

?Como se comporta una regresion basada en SVM para estimar serenazgo? &

Mismas variables de regresion que MLP para mantener comparabilidad. &

Se implementa SVR lineal con perdida epsilon-insensitiva. Tiene resultados positivos, pero mayor RMSE que MLP. Codigo: \texttt{SVRegressor.scala}. \\

\midrule

Clasificacion Random Forest &

¿Se puede clasificar un predio como exclusivo o compartido usando variables decimales sin fuga de datos? &

\texttt{TIPO\_PREDIO} como etiqueta derivada; se excluye \texttt{PORCENTAJE\_CONDOMINIO} de las features para evitar fuga. &

Random Forest logra el mejor F1-score. Su interpretacion es que los arboles capturan mejor reglas por umbral y relaciones no lineales. Codigo: \texttt{RFAndMLPClassification.scala}. \\

\midrule

Clasificacion MLP &

¿Que desempeno tiene una red neuronal para clasificar tipo de predio frente a Random Forest? &

Mismas variables de clasificacion que Random Forest. &

MLPClassifier funciona, pero obtiene menor F1-score que Random Forest, por lo que requiere mayor ajuste. Codigo: \texttt{RFAndMLPClassification.scala}. \\

\bottomrule

\end{longtable}